\section{جمع‌بندی و پیشنهادها}

در این پژوهش، چارچوبی مبتنی بر حسگرهای اینرسیایی پوشیدنی و روش‌های یادگیری ماشین برای ثبت، پردازش و طبقه‌بندی الگوهای حرکتی متناظر با سطوح عملکردی تعریف‌شده در تعدادی از فعالیت‌های مقیاس \lr{NSAA} طراحی و پیاده‌سازی شد. هدف اصلی، بررسی امکان تفکیک سه سطح عملکردی شبیه‌سازی‌شده با استفاده از اطلاعات شتاب خطی و سرعت زاویه‌ای بخش‌های مختلف بدن بود.

برای دستیابی به این هدف، ابتدا سامانه‌ای شامل پنج حسگر اینرسیایی و یک نرم‌افزار اختصاصی برای دریافت، نمایش و ذخیره‌سازی هم‌زمان داده‌ها توسعه یافت. سپس داده‌های ثبت‌شده پس از اصلاح خطای پیش‌مقدار، فیلترگذاری و آماده‌سازی، وارد مرحله استخراج ویژگی شدند. از هر کانال سیگنال مجموعه‌ای از ویژگی‌های آماری و حوزه زمان استخراج شد و در ادامه، اثر کاهش بعد به روش تحلیل مؤلفه‌های اصلی و عملکرد سه الگوریتم ماشین بردار پشتیبان، نزدیک‌ترین همسایه و جنگل تصادفی مورد بررسی قرار گرفت.

لازم است تأکید شود که داده‌های این پژوهش از بیماران مبتلا به دیستروفی عضلانی دوشن ثبت نشده‌اند. داده‌ها توسط یک فرد سالم و آشنا با بیماری و معیار امتیازدهی \lr{NSAA} و با هدف شبیه‌سازی کنترل‌شده درجات عملکردی تولید شده‌اند. بنابراین، نتایج این پژوهش نشان‌دهنده امکان‌پذیری فنی روش پیشنهادی در داده‌های شبیه‌سازی‌شده است و نباید به‌عنوان اعتبارسنجی بالینی سامانه یا معادل مستقیم عملکرد بیماران تفسیر شود.

\subsection{دستاوردهای اصلی}

نخستین دستاورد پژوهش، طراحی یک زنجیره منسجم برای ثبت و مدیریت داده‌های حرکتی بود. در این زنجیره، اطلاعات شتاب خطی و سرعت زاویه‌ای پنج حسگر به‌صورت هم‌زمان دریافت و همراه با مشخصات فعالیت، تکرار آزمون و سطح عملکردی شبیه‌سازی‌شده ذخیره شدند. این ساختار امکان ایجاد مجموعه‌داده‌ای منظم و قابل استفاده در مراحل بعدی پردازش و مدل‌سازی را فراهم کرد.

دستاورد دوم، پیاده‌سازی کامل مسیر پردازش از سیگنال خام تا ورودی مدل‌های یادگیری ماشین بود. این مسیر شامل بررسی کیفیت داده‌ها، اصلاح خطای پیش‌مقدار ژیروسکوپ، اعمال فیلتر پایین‌گذر باترورث، بررسی اثر جداسازی بازه فعالیت، استخراج ویژگی، استانداردسازی و کاهش بعد بود. قرار دادن استانداردسازی و تحلیل مؤلفه‌های اصلی درون \lr{Pipeline} نیز موجب شد این تبدیل‌ها فقط بر داده‌های آموزشی برازش شوند و از نشت اطلاعات میان مجموعه‌های آموزش و آزمون جلوگیری شود.

\subsection{جمع‌بندی نتایج فعالیت‌ها}

خلاصه تنظیمات مدل‌های منتخب برای پنج فعالیت بررسی‌شده در جدول \ref{tab:final_selected_models} ارائه شده است. حد نصاب‌های ذکرشده، معیار داخلی مورد استفاده برای انتخاب تعداد مؤلفه‌ها بوده‌اند و نباید به‌عنوان دقت بالینی سامانه تفسیر شوند.

\begin{table}[htbp]
	\centering
	\caption{خلاصه مدل‌های منتخب برای فعالیت‌های بررسی‌شده}
	\label{tab:final_selected_models}
	\begin{tabular}{c c c c c}
		\hline
		فعالیت & مدل منتخب & تعداد مؤلفه اصلی & حد نصاب دقت & جداسازی فعالیت \\
		\hline
		ایستادن & \lr{SVM} & ۱۳ & ۹۰ درصد & بدون جداسازی \\
		برخاستن از صندلی & \lr{SVM} & ۱۴ & ۸۰ درصد & بدون جداسازی \\
		پریدن & \lr{SVM} & ۱۵ & ۹۰ درصد & بدون جداسازی \\
		بالا آوردن سر & \lr{SVM} & ۱۴ & ۹۰ درصد & بدون جداسازی \\
		بالا رفتن از جعبه با پای راست & \lr{SVM} & ۱۰ & ۸۰ درصد & بدون جداسازی \\
		\hline
	\end{tabular}
\end{table}

در فعالیت ایستادن، دامنه تغییرات حرکتی نسبت به فعالیت‌های دینامیک محدودتر بود و تفاوت سطوح عملکردی بیشتر در نوسان‌های کوچک بدن، نحوه حفظ تعادل و میزان حرکت اندام‌های فوقانی ظاهر شد. با وجود ماهیت کم‌دامنه این فعالیت، مدل ماشین بردار پشتیبان با ۱۳ مؤلفه اصلی توانست حد نصاب ۹۰ درصد را تأمین کند. این نتیجه نشان می‌دهد ویژگی‌های آماری استخراج‌شده می‌توانند بخشی از تفاوت‌های ظریف میان الگوهای شبیه‌سازی‌شده ایستادن را بازنمایی کنند.

فعالیت برخاستن از صندلی شامل انتقال وزن، حرکت تنه، بازشدن مفاصل اندام تحتانی و احتمال استفاده جبرانی از دست‌ها بود. به همین دلیل، اطلاعات ثبت‌شده از حسگرهای سر، دست‌ها و پاها در کنار یکدیگر برای تفکیک الگوهای اجرا اهمیت داشتند. مدل \lr{SVM} با ۱۴ مؤلفه اصلی به نزدیک‌ترین عملکرد پایدار نسبت به حد نصاب ۸۰ درصد رسید. نتیجه این فعالیت نشان داد که تفاوت راهبردهای حرکتی در یک حرکت انتقالی چندمرحله‌ای در ویژگی‌های استخراج‌شده قابل مشاهده است، هرچند عملکرد آن نسبت به برخی فعالیت‌های دیگر محدودتر بود.

در فعالیت پریدن، تغییرات سریع شتاب و سرعت زاویه‌ای، جداشدن پاها از زمین و ضربه ناشی از فرود، الگوهای مشخص‌تری در سیگنال‌ها ایجاد کردند. مدل ماشین بردار پشتیبان با ۱۵ مؤلفه اصلی حد نصاب ۹۰ درصد را تأمین کرد. در این فعالیت، ویژگی‌های مرتبط با دامنه تغییرات، انرژی سیگنال، پراکندگی و نوسان‌های سریع نقش مهمی در نمایش تفاوت درجات عملکردی شبیه‌سازی‌شده داشتند.

در فعالیت بالا آوردن سر، حسگر نصب‌شده روی پیشانی نقش اصلی را در ثبت دامنه و سرعت حرکت سر ایفا کرد، درحالی‌که حسگرهای دست و پا اطلاعات مکملی درباره حرکات جبرانی احتمالی فراهم کردند. مدل \lr{SVM} با ۱۴ مؤلفه اصلی به حد نصاب ۹۰ درصد رسید. این نتیجه بیانگر آن است که ترکیب ویژگی‌های مربوط به حرکت اصلی سر و تغییرات سایر اندام‌ها می‌تواند برای تفکیک الگوهای شبیه‌سازی‌شده این فعالیت مفید باشد.

در فعالیت بالا رفتن از جعبه با پای راست، تغییرات حسگر پای راست برای توصیف آغاز حرکت و انتقال وزن و حسگر پای چپ برای ثبت تکمیل حرکت اهمیت بیشتری داشتند. اطلاعات حسگر پیشانی و دست‌ها نیز نحوه کنترل تنه و حفظ تعادل را منعکس کردند. برای این فعالیت، ماشین بردار پشتیبان با ۱۰ مؤلفه اصلی حد نصاب ۸۰ درصد را تأمین کرد. کمتر بودن تعداد مؤلفه‌های موردنیاز نشان می‌دهد بخش مهمی از اطلاعات تفکیک‌کننده این فعالیت در فضای نسبتاً فشرده‌ای از مؤلفه‌ها قرار داشته است.

بررسی جداسازی زمانی فعالیت‌ها نشان داد که در هیچ‌یک از پنج فعالیت، استفاده از سیگنال جداسازی‌شده به بهبود پایدار عملکرد مدل نهایی منجر نشد. به همین دلیل، مدل‌های نهایی بر داده‌های بدون جداسازی آموزش داده شدند. این نتیجه به معنای غیرضروری بودن جداسازی فعالیت در کاربردهای آینده نیست؛ زیرا ممکن است بخش‌های پیش و پس از حرکت در داده‌های حاضر دارای اطلاعاتی درباره وضعیت اولیه، آماده‌سازی یا بازیابی تعادل بوده باشند. اثر جداسازی باید در مجموعه‌داده‌های بزرگ‌تر و با الگوریتم‌های دقیق‌تر دوباره بررسی شود.

در مجموع، نتایج پنج فعالیت نشان دادند که ماهیت حرکت بر تعداد مؤلفه‌های مناسب، میزان اطلاعات موجود در سیگنال‌ها و سطح عملکرد مدل اثر می‌گذارد. در عین حال، انتخاب ماشین بردار پشتیبان برای همه فعالیت‌ها نشان داد این مدل در فضای ویژگی کاهش‌یافته، برای تفکیک درجات عملکردی شبیه‌سازی‌شده از پایداری مناسبی برخوردار بوده است.

\subsection{محدودیت‌های پژوهش}

مهم‌ترین محدودیت پژوهش، استفاده از داده‌های شبیه‌سازی‌شده توسط یک فرد سالم است. الگوهای حرکتی بیماران مبتلا به دیستروفی عضلانی دوشن ممکن است تحت تأثیر شدت بیماری، سن، خستگی، کوتاهی عضلات، محدودیت دامنه حرکتی و راهبردهای جبرانی فردی قرار گیرند و پیچیدگی بیشتری نسبت به الگوهای کنترل‌شده این پژوهش داشته باشند. ازاین‌رو، قابلیت مدل‌ها در تفکیک داده‌های حاضر را نمی‌توان مستقیماً به بیماران تعمیم داد.

از آنجا که تمام داده‌ها توسط یک فرد تولید شده‌اند، اعتبارسنجی متقاطع انجام‌شده عمدتاً پایداری مدل در برابر تکرارهای مختلف همان فرد را ارزیابی می‌کند و توانایی تعمیم مدل به افراد جدید را نشان نمی‌دهد. ارزیابی واقعی تعمیم‌پذیری مستلزم ثبت داده از چندین فرد و استفاده از روش‌هایی مانند اعتبارسنجی مستقل از آزمودنی است.

تعداد نمونه‌ها نیز در مقایسه با تعداد ویژگی‌های اولیه محدود بود. کاهش بعد و استفاده از اعتبارسنجی متقاطع، خطر بیش‌برازش را کاهش می‌دهند، اما آن را به‌طور کامل حذف نمی‌کنند. همچنین انتخاب مدل‌ها بر اساس همین مجموعه‌داده انجام شده است؛ بنابراین، ارزیابی نهایی باید روی یک مجموعه آزمون کاملاً مستقل صورت گیرد.

نتایج به نحوه نصب حسگرها، جهت‌گیری آن‌ها، میزان سفتی اتصال به بدن و شرایط اجرای آزمون وابسته‌اند. تغییر در محل یا زاویه حسگر می‌تواند توزیع سیگنال‌ها و مقادیر ویژگی‌ها را تغییر دهد. افزون بر این، کاهش نرخ نمونه‌برداری در برخی بازه‌های استفاده طولانی از حسگرها نشان می‌دهد که پایداری سخت‌افزار و منبع تغذیه باید پیش از استفاده گسترده‌تر بهبود یابد.

ویژگی‌های استفاده‌شده عمدتاً آماری و مربوط به حوزه زمان بودند. این ویژگی‌ها برای توسعه اولیه سامانه مناسب و قابل تفسیرند، اما ممکن است تمام اطلاعات موجود در ساختار زمانی و فرکانسی حرکت را استخراج نکنند. علاوه بر این، تنها سه الگوریتم کلاسیک با تنظیمات مشخص بررسی شدند و بهینه‌سازی گسترده ابرپارامترها انجام نشد.

در نهایت، حد نصاب‌های ۸۰ و ۹۰ درصد به‌عنوان معیار مهندسی برای انتخاب مدل در مجموعه‌داده حاضر در نظر گرفته شدند. عبور از این حد نصاب‌ها به‌تنهایی بیانگر روایی، پایایی یا کاربردپذیری بالینی سامانه نیست.

\subsection{پیشنهادها برای کارهای آینده}

مهم‌ترین گام آینده، جمع‌آوری داده از بیماران مبتلا به دیستروفی عضلانی دوشن با همکاری تیم بالینی است. اجرای هر فعالیت باید هم‌زمان توسط ارزیابان متخصص بر اساس دستورالعمل رسمی \lr{NSAA} امتیازدهی شود تا امکان مقایسه خروجی سامانه با نمره بالینی فراهم گردد. حضور چند ارزیاب می‌تواند برای بررسی توافق بین ارزیابان و مقایسه آن با مدل نیز مفید باشد.

داده‌برداری باید از افراد متعدد، در جلسات مختلف و در سطوح متفاوت بیماری انجام شود. در این شرایط استفاده از اعتبارسنجی مستقل از آزمودنی، مانند کنار گذاشتن کامل داده‌های یک فرد در هر مرحله آزمون، ضروری است. همچنین بهتر است یک مجموعه داده مستقل برای ارزیابی نهایی مدل پس از پایان انتخاب ویژگی و تنظیم ابرپارامترها نگه داشته شود.

گسترش فعالیت‌های بررسی‌شده به سایر اجزای مقیاس \lr{NSAA} می‌تواند امکان تخمین بخش بزرگ‌تری از عملکرد حرکتی را فراهم کند. بااین‌حال، پیش از افزودن همه فعالیت‌ها لازم است قابلیت اندازه‌گیری هر فعالیت با حسگرهای اینرسیایی و نیاز احتمالی آن به تجهیزات مکمل بررسی شود.

بهبود الگوریتم تشخیص شروع و پایان حرکت، هم‌زمان‌سازی دقیق‌تر حسگرها، تثبیت نرخ نمونه‌برداری و تدوین دستورالعمل استاندارد نصب حسگرها از دیگر اقدامات ضروری است. مقایسه داده‌های کامل، بازه‌های جداسازی‌شده و پنجره‌های زمانی مختلف می‌تواند مشخص کند کدام بخش سیگنال بیشترین اطلاعات را برای هر فعالیت فراهم می‌کند.

در بخش مدل‌سازی، پیشنهاد می‌شود انتخاب ویژگی و کاهش بعد به‌صورت مستقل برای هر فعالیت بهینه‌سازی شود. استفاده از اعتبارسنجی تو‌در‌تو برای تنظیم ابرپارامترهای مدل، بررسی ویژگی‌های حوزه فرکانس و زمان--فرکانس و مقایسه روش‌های انتخاب ویژگی با تحلیل مؤلفه‌های اصلی می‌تواند به افزایش قابلیت تعمیم مدل کمک کند. در صورت افزایش حجم داده‌ها، استفاده از مدل‌های مبتنی بر توالی زمانی و روش‌های یادگیری عمیق نیز قابل بررسی خواهد بود.

همچنین می‌توان با حذف مرحله‌ای حسگرها، سهم هر محل نصب را در عملکرد مدل بررسی کرد. چنین تحلیلی مشخص می‌کند آیا امکان کاهش تعداد حسگرها و ساده‌تر شدن سامانه، بدون افت قابل توجه عملکرد، وجود دارد یا خیر.

در نهایت، توسعه یک رابط نرم‌افزاری یکپارچه برای دریافت داده، کنترل کیفیت، پردازش، اجرای مدل و تولید گزارش می‌تواند سامانه را برای استفاده در مطالعات بالینی آینده آماده‌تر کند. این رابط باید علاوه بر نمره پیش‌بینی‌شده، میزان اطمینان مدل و شاخص‌های قابل تفسیر مربوط به نحوه اجرای حرکت را نیز نمایش دهد.

\subsection{نتیجه‌گیری نهایی}

در این پژوهش، یک سامانه کامل از مرحله ثبت داده‌های اینرسیایی تا طبقه‌بندی درجات عملکردی شبیه‌سازی‌شده طراحی و آزمایش شد. نتایج نشان دادند که ویژگی‌های استخراج‌شده از شتاب خطی و سرعت زاویه‌ای حسگرهای نصب‌شده روی سر، دست‌ها و پاها، در مجموعه‌داده حاضر دارای اطلاعات کافی برای تفکیک سه سطح شبیه‌سازی‌شده پنج فعالیت منتخب بوده‌اند.

در مدل‌های نهایی، ماشین بردار پشتیبان همراه با تحلیل مؤلفه‌های اصلی برای همه فعالیت‌ها انتخاب شد، اما تعداد مناسب مؤلفه‌ها با توجه به ماهیت هر حرکت متفاوت بود. دستیابی به عملکرد موردنظر با ۱۰ تا ۱۵ مؤلفه اصلی نشان داد می‌توان بخش عمده اطلاعات موردنیاز برای طبقه‌بندی را در فضایی بسیار کوچک‌تر از بردار اولیه ۶۹۰ ویژگی نمایش داد.

بااین‌حال، دستاورد اصلی پژوهش اثبات کاربرد بالینی سامانه نیست، بلکه ارائه یک چارچوب فنی و روش‌مند برای ثبت، پردازش و تحلیل کمی حرکات مرتبط با \lr{NSAA} است. اعتبار واقعی این چارچوب باید از طریق داده‌برداری از بیماران، مقایسه با ارزیابی متخصصان و آزمون روی افراد و جلسات مستقل بررسی شود.

در صورت تکمیل این مراحل، سامانه پیشنهادی می‌تواند در آینده به‌عنوان ابزاری مکمل برای عینی‌تر و تکرارپذیرتر شدن ارزیابی عملکرد حرکتی، بررسی تغییرات بیمار در طول زمان و توسعه روش‌های پایش از راه دور مورد استفاده قرار گیرد.